\documentclass[a4paper,12pt]{article}
\usepackage[slovene]{babel}
\usepackage[utf8]{inputenc}
\usepackage[T1]{fontenc}
\usepackage{lmodern}

\usepackage{amsmath,amsfonts}

\setlength{\parindent}{0mm}

\begin{document}

% Datoteko dopolnite tako, da bo PDF čim bolj podoben datoteki
% matematicni-izrazi-resitev.pdf.
% Kjer piše ??, manjka matematični izraz.
% Matematičnih nalog ne rešujte, samo zapišite matematične izraze.

\section*{Logika in množice}

\begin{enumerate}
\item
Poišči preneksno obliko formule ??.

\item
Definiramo množici ?? in ??.
V ravnino nariši:
\begin{enumerate}
\item ??
\item ??
\end{enumerate}

\item
Dokaži:
\begin{itemize}
\item ??
\item ??
\end{itemize}
\end{enumerate}

\section*{Analiza}

\begin{enumerate}
\item
Pokaži, da je funkcija ?? enakomerno zvezna na ??.

\item
Katero krivuljo določa naslednji parametrični zapis?
% Pomagajte si z dokumentacijo o razmikih v matematičnem načinu:
% https://www.overleaf.com/learn/latex/Spacing_in_math_mode
\[
x(t) = a \cos t, ?? % tu manjka ukaz za presledek
y(t) = b \sin t, ?? % tu manjka ukaz za presledek
t \in [0, 2 \pi]
\]

\item
Pokaži, da ima ?? inverzno funkcijo in izračunaj ??.

\item
Izračunaj integral
% V rešitvi je integral zapisan v vrstičnem načinu, a v prikaznem slogu.
% To dosežete tako, da v matematičnem načinu najprej uporabite ukaz displaystyle.
% Pred dx je presledek: pravi ukaz je \,
??
% \frac{2+\sqrt{x+1}}{(x+1)^2-\sqrt{x+1}}

\item
Naj bo \(g\) zvezna funkcija. Ali posplošeni integral
??
konvergira ali divergira? Utemelji.

\item
S Taylorjevo formulo pokaži, da v okolici točke \(0\) velja
??
\end{enumerate}

\section*{Kompleksna števila}

\begin{enumerate}
\item
Naj bo \(z\) kompleksno število, \(z \ne 1\) in ??.
Dokaži, da je število \( i \, \frac{z+1}{z-1} \) realno.

\item
Poenostavi izraz:
% V rešitvi so notranji ulomki zapisani z ukazom dfrac.
??
\end{enumerate}

\section*{Zaporedja, vrste in limite}

\begin{enumerate}
\item
Z indukcijo dokaži, da za vsako naravno število \(n\) velja
??

\item
Naj bo ?? absolutno konvergentna vrsta in \(a_n \ne -1\).
Dokaži, da je tudi vrsta \(\sum_{n=1}^\infty \frac{a_n}{1+a_n}\)
absolutno konvergentna.

\item
Izračunaj limito
??

\item
Za dani zaporedji preveri, ali sta konvergentni.
% Pomagajte si s spodnjima delno pripravljenima matematičnima izrazoma:
% a_n = \sqrt{2+\sqrt{2+\dots+\sqrt{2}}} \qquad
% b_n = \sin(\sin(\dots(\sin 1)\dots))
% Zaviti oklepaj pod izrazom naredi ukaz underbrace: \underbrace{izraz}_{napis}.
% Besedilo v napisu napišite z ukazom text, presledek pred njim pa z znakom ~.
??
\end{enumerate}

\section*{Algebra}

\begin{enumerate}
\item
Vektorja ??
sta pravokotna in imata dolžino 1. Določi kot med vektorjema \(\vec{a}\) in \(\vec{b}\).

\item
Izračunaj
??
% 1 2 3 4 5 6
% 4 5 2 6 3 1
\end{enumerate}

\section*{Velika determinanta}

Izračunaj naslednjo determinanto \(2n \times 2n\), ki ima na neoznačenih mestih ničle.
% Prazna mesta v determinanti dobite tako, da med znakoma & ne napišete ničesar.
% Pike naredijo ukazi cdots (vodoravne), vdots (navpične) in ddots (poševne).
??
% 1  1
% 2  2
%
% n-1  n-1
% 1  2  n-1  n  n+1  n+2  2n
% n+1  n+1
% n+2  n+2
%
% 2n  2n

\section*{Grupe}

Naj bo
% Definiciji množic G in H sta v rešitvi poravnani pri enačaju, z okoljem align*.
% Zavita oklepaja napišete kot \{ in \}, presledek za podpičjem pa z ukazom \;
??
\begin{enumerate}
\item
Pokaži, da je \(G\) podgrupa v grupi ??
neničelnih kompleksnih števil za običajno množenje.

\item
Pokaži, da je \(H\) podgrupa v aditivni grupi ??
ravninskih vektorjev za običajno seštevanje po komponentah.

\item
Pokaži, da je preslikava \(f:H\to G\), podana s pravilom
??
izomorfizem grup \(G\) in \(H\).
\end{enumerate}

\end{document}
